\documentclass[10pt,a4paper]{amsart}
\usepackage[margin=19mm]{geometry}
\usepackage{amsmath,amssymb,mathtools}
\usepackage[scr=rsfs,cal=euler]{mathalfa}
\usepackage{xfrac}
\usepackage[unicode,hidelinks,bookmarks=false]{hyperref}
\newcommand{\ie}{i.e.\ }
\newcommand{\cf}{cf.\ }
\newcommand{\resp}{respectively\ }
\newcommand{\Subsection}{\subsection}
\providecommand{\nobreakdash}{\nobreak\hbox{-}}
\setlength{\emergencystretch}{2em}
\multlinegap0pt
\usepackage{amssymb}
\newtheorem{thm}{Theorem}
 \newcommand{\RM}{\mathbb{R}}
 \newcommand{\dlambda}{\operatorname{d\lambda}} 
 \newcommand{\ds}{\operatorname{ds}}
\title{Fourier inversion theorems for integral transforms involving Bessel functions}
\author{Alexey Gorshkov}
\date{Source: arXiv:2111.12674v2 (CC BY 4.0)}
\begin{document}
\maketitle

\noindent\emph{Source and licence.} Fourier inversion theorems for integral transforms involving Bessel functions. Version arXiv:2111.12674v2 (CC BY 4.0), distributed by arXiv under the Creative Commons Attribution 4.0 International licence, http://creativecommons.org/licenses/by/4.0/, as recorded in the arXiv licence metadata for this exact version and shown on the abstract page https://arxiv.org/abs/2111.12674v2. Full licence notice: \texttt{licenses/arxiv-2111.12674-CC-BY-4.0.txt}.\par\smallskip

\noindent\textit{Editorial scope.} Counted unit: the article's thm wkk:thm (source0.tex lines 148--156). The article's complete original proof of it is delivered with it (source0.tex lines 772--787, one \texttt{proof} environment of 16 lines, which the article places away from the statement). No mathematics is written, repaired or re-derived: the statement and the proof are the article's own lines, in the article's own notation, and no retained line is substituted or re-flowed.  Retained uncounted as the source-local context the proof needs: lines 116--120; lines 389--389; lines 473--475; lines 518--534; lines 547--549; lines 760--766.  Nothing was omitted from any retained span, so the package carries no omission record.  No retained line contains a citation, so the wrapper carries no bibliography and no bibliography entry.  No retained line contains a figure environment, an image-inclusion command or drawing code, so no image is delivered and the package declares no asset dependency.  Editorial formatting only, in addition: an AMS \texttt{amsart} preamble that restates the article's own declarations and macros used by the retained lines, namely the environments \texttt{thm} on separate counters, together with the standard packages those lines need.  The retained environments print in this document as Theorem 1 in order of appearance, whereas the article prints the same items with the numbering of its own full document.  The complete native source of the article as distributed in the arXiv e-print for this exact version is bundled unchanged as \texttt{sources/118853/source0.tex}.
\begin{align}\label{int:weberorr}
W_{k,l}[f](\lambda) = \int_{r_0}^\infty \frac{J_{k}(\lambda s)Y_{l}(\lambda r_0) - Y_{k}(\lambda s)J_{l}(\lambda r_0)}{\sqrt{J_{l}^2(\lambda r_0) + Y_{l}^2(\lambda r_0)}} f(s) s \ds,k,l\in \RM,
\\ \label{int:weberorrinv}
W^{-1}_{k,l}[\hat f](r) = \int_{0}^\infty \frac{J_{k}(\lambda r)Y_{l}(\lambda r_0) - Y_{k}(\lambda r)J_{l}(\lambda r_0)}{\sqrt{J_{l}^2(\lambda r_0) + Y_{l}^2(\lambda r_0)}} \hat f (\lambda) \lambda \dlambda.
\end{align}

\subsection{Proof of the inversion theorem} \label{wkk:proof}

\begin{equation} \label{wkk1:omegaeqpolar}
\frac{\partial w(t,r)}{\partial t} - \Delta_k w(t,r) = 0,
\end{equation}

\begin{thm}\label{wkk1:thm}Let $f(r) \sqrt r \in L_1(r_0,\infty)\cap L_2(r_0,\infty)$, $r_0>0$, $k \in \RM$. Then the associated Weber-Orr transforms (\ref{int:weberorr}), (\ref{int:weberorrinv}) satisfy  almost everywhere 
\begin{align*}
f(r) &= W^{-1}_{k,k-1}\left [W_{k,k-1} [f] \right ](r), k\leq 1, \\
f(r) &= W^{-1}_{k,k-1}\left [W_{k,k-1} [f] \right ](r) +  \frac {2(k-1) r_0^{2(k-1)}} {r^k} \int_{r_0}^\infty s^{-k+1} f(s) \ds,k>1,\\
f(r) &= W^{-1}_{k,k+1}\left [W_{k,k+1} [f] \right ](r), k\geq -1, \\
f(r) &= W^{-1}_{k,k+1}\left [W_{k,k+1} [f] \right ](r) -  \frac {2(k+1) r^{k}} {r_0^{2(k+1)}} \int_{r_0}^\infty s^{-k+1} f(s) \ds,k<-1.
\end{align*}
and the following Plancherel-Parseval equity holds:
\begin{align*}
\|f\|_{L_2(r_0,\infty; r)}^2 &= \| W_{k,k-1}[f] \|_{L_2(0,\infty; r) }^2 , k\leq 1, \\
\|f\|_{L_2(0,\infty; r) }^2 &= \| W_{k,k-1}[f] \|_{L_2(0,\infty; r) } ^2 \\
&+ {2(k-1)} r_0^{2(k-1)} \left (f,1/r^k \right )^2_{L_2(r_0,\infty; r)}, k> 1,\\
\|f\|_{L_2(r_0,\infty; r)}^2 &= \| W_{k,k+1}[f] \|_{L_2(0,\infty; r) }^2 , k\geq -1, \\
\|f\|_{L_2(0,\infty; r) }^2 &= \| W_{k,k-1}[f] \|_{L_2(0,\infty; r) } ^2 \\
&- \frac{2(k+1)}{r_0^{2(k+1)}} \left (f,r^k \right )^2_{L_2(r_0,\infty; r)}, k< -1.
\end{align*}
\end{thm}

\begin{equation}\label{wkk1:robin_bound1}
r_0\frac{\partial w(t,r)}{\partial r}\Big|_{r=r_0} + k w(t,r_0) = 0,~k \in \RM.
\end{equation}

\begin{align}
{\rm for~} k \leq 1:\nonumber  \\
&w(t,r) = W^{-1}_{k,k-1}\left [ e^{-\lambda^2 t} W_{k,k-1}[f] \right ], \label{wkk1:heatsol1}\\
{\rm for~} k > 1: \nonumber \\
&w(t,r) = W^{-1}_{k,k-1}\left [ e^{-\lambda^2 t} W_{k,k-1}[f] \right ]\label{wkk1:heatsol2}\\
&+ \frac {2(k-1) r_0^{2k-2}} {r^{k}} \left (\frac 1 r^{k}, f \right )_{L_2(r_0,\infty; r)} \nonumber
\end{align}

\begin{thm}\label{wkk:thm}Let $f(r) \sqrt r \in L_1(r_0,\infty)\cap L_2(r_0,\infty)$, $k \in \RM$, $r_0>0$. Then the  Weber-Orr transforms $W_{k,k}[\cdot]$ defined by (\ref{int:weberorr}), (\ref{int:weberorrinv}) satisfy  almost everywhere 
\begin{align}\label{wkk:inversion_formula}
f(r) = W^{-1}_{k,k}\left [W_{k,k} [f] \right ](r), 
\end{align}
and the following Plansherel equity holds:
$$
||f||_{L_2(r_0,\infty; r)} = \| W_{k,k}[f] \|_{L_2(0,\infty; r) } 
$$
\end{thm}

\begin{proof}

For $k\leq 1$ the proof of the theorem \ref{wkk1:thm} repeats word by word the proof of theorem \ref{wkk:thm} in section \ref{wkk:proof}. We must prove the validity of the limit transition $t\to 0$ in (\ref{wkk1:heatsol1}). Let $w(t,r)$ be solution of (\ref{wkk1:omegaeqpolar}), (\ref{wkk1:robin_bound1}) with initial datum $f(r)$. $\Delta_k$ generates strongly continuous semi-group in $L_2(r_0,\infty; r)$ and 
$$
\|w(t,\cdot) - f(\cdot)\| \to 0,~t\to 0.
$$ 

Then as in section \ref{wkk:proof} one could get
$$
w(t,\cdot) \rightharpoondown W^{-1}_{k,k-1}\left [W_{k,k-1} [f] \right ],~t\to 0
$$
and so almost everywhere holds
$$f(\cdot)=W^{-1}_{k,k-1}\left [W_{k,k-1} [f] \right ]\|.$$

For $k>1$ we have addition stationary term in (\ref{wkk1:heatsol2}) and also can pass to limit $t\to 0$ obtaining the required inversion formula.
\end{proof}

\end{document}
